\thispagestyle{plain}

% Title with Titles Font (Anonymous Pro)
{\titlesfont\fontsize{18}{22}\textbf{\color{cherenkovblue}{Decarbonizing Sette Laghi \\ Technical, Economic, and Environmental Study}}}\\
{\titlesfont\fontsize{11}{13}\color{cherenkovblue} Politecnico di Milano}\\

\vspace{-8pt}

% Author information with Main Font (Poppins)
{\normalsize\textbf{Alessandro Borgonovo, Gaia Larici, Riccardo Monti, Simone Pagliuca, Riccardo Ronchi}} \\
{\footnotesize\textit{borgo, gaia, riccardo, 10698534, richard}} \\

\vspace{8pt}

% Course and Academic Year
{\footnotesize\textbf{Course:} Integration of Nuclear and Renewables for Carbon Neutral Scenarios}\\
{\footnotesize\textbf{Academic Year:} 2024/2025}\\

% Horizontal rule
\vspace{6pt}
\centerline{\rule{0.95\textwidth}{0.5pt}}
\vspace{8pt}

% Abstract with Titles Font for Heading, Main Font for Body
{\titlesfont\fontsize{10}{12}\textbf{\color{cherenkovblue} ABSTRACT}} 

\vspace{4pt}
{\normalsize
This study evaluates the technical and economic feasibility of integrating Small Modular Reactors (SMRs) into the residential energy sector of the north-western province of Varese, Italy. The proposed hybrid system combines nuclear power with renewable energy sources, hydrogen production, and synthetic methane generation to meet regional energy demands and support decarbonization efforts. The analysis includes dynamic simulations of the plant's operation, economic assessments, and environmental impact evaluations over a 60-year period. The methodology involves data collection from regional energy databases, dynamic modeling using Dymola software, and economic analysis considering capital and operational expenditures, as well as revenue streams from electricity, hydrogen, and synthetic gas sales. Preliminary results indicate the potential for significant decarbonization and energy stability, although further analysis is required to determine the economic viability and optimal configuration of the system. [Placeholder for conclusions: The study concludes with insights into the overall feasibility and potential improvements for the proposed hybrid energy system.]
}

\vspace{10pt}

% Key-words box with Titles Font
\begin{tcolorbox}[arc=0pt, boxrule=0pt, colback=cherenkovblue!60, width=\textwidth, colupper=white]
    {\titlesfont\fontsize{10}{12}\textbf{Key-words:}} Small Modular Reactors (SMRs), Hybrid Energy Systems, Decarbonization, Hydrogen Production, Synthetic Methane, Dynamic Simulation, Economic Feasibility, Environmental Impact, Residential Energy Sector, Renewable Energy Integration
\end{tcolorbox}

\vspace{12pt}

% Nomenclature Definitions
\makenomenclature
\renewcommand\nomgroup[1]{%
\item[\bfseries
\ifstrequal{#1}{A}{Type A}{%
\ifstrequal{#1}{B}{Type B \& Misc}{%
}}%
]}

% Neutronics Quantities
\nomenclature[A]{$math$}{Description ($units-of-measurement$)}

% Data & Statistics
\nomenclature[B]{$math$}{Description ($units-of-measurement$)}

% Two-column layout for the nomenclature
\renewcommand{\nompreamble}{\begin{multicols}{2}}
\renewcommand{\nompostamble}{\end{multicols}}

% Render Nomenclature Without Section Break
\vspace{-5pt} % Adjust spacing to fit table
\printnomenclature
